\section{Illustrative Stakeholder-Specific Narratives}
\label{SEC:stakeholder_narrative_examples}
% =========================================================

To demonstrate stakeholder-specific evidence selection, six stakeholder cards were generated from
the same structured evidence. Each card answers its stakeholder's primary question with panels
drawn by the dashboard itself, four key figures and a short two-part narrative. Rather than
forcing every role through one case, each card selects the instance that role needs: the resident
card explains a typical area, Stoneferry (Kingston upon Hull 016B, \texttt{E01012806}); the police
card examines the city-centre and Old Town LSOA (029E, \texttt{E01033104}); the manager card
contrasts the Hull-wide pattern with Bilton Grange (010B, \texttt{E01012825}), the most geographically sensitive of the areas not led by crime history; and the planner,
researcher and governance cards report results across Hull.

All values come from the final Hull design: the Feature branch explains the 30-predictor
XGBoost model, and the random-donor Geo branch explains an XGBoost model on the same predictors
plus the centroid coordinates, in which masking an LSOA moves its coordinates to the centroids of
randomly drawn other Hull LSOAs ($N=4{,}000$ perturbations, 50 donor draws, seed 42). Place names and streets
come from the \texttt{place\_context} supplied to the prompt, built from OpenStreetMap features
inside the official 2021 LSOA boundaries.

The text, key figures and highlighted figures and terms of every card were generated by GPT-4.1
through the dashboard's LLM pipeline at temperature 0, from the master prompt
(\texttt{geosmile-hull-v9}) followed by a stakeholder-specific card-mode addendum; none was written
or edited by hand. Each output was checked automatically before use: every cited figure must be
traceable to the evidence package, each section must respect its length and key-term limits, and
the cards for non-technical roles may not use test statistics. Where a run failed these checks,
the addendum was revised and the card regenerated; the prompts, evidence, raw output and checks
of every run are archived. The cards correspond to the dashboard's short mode; the dashboard also
provides a longer, comprehensive narrative for each role.

The examples are intended to demonstrate differences in evidence selection, rather than merely
differences in vocabulary. Each subsection therefore introduces the stakeholder's questions and
presents the card on the same page.
